% Metric schema interactive_nav_v3_paper_metrics_v4. Requires amsmath.
\paragraph{Evaluation metrics.}
Let $S_i$ indicate verified navigation success, $L_i$ the executed planar
base path length, and $L_i^{\mathrm{ref}}$ the frozen feasible reference
length. We report episode means of SR and SPL:
\begin{equation}
\mathrm{SR}=\frac{1}{N}\sum_i S_i,\qquad
\mathrm{SPL}=\frac{1}{N}\sum_i
 S_i\frac{L_i^{\mathrm{ref}}}{\max(L_i^{\mathrm{ref}},L_i)}.
\end{equation}
A successful zero-reference, zero-motion episode has SPL one. Missing or
invalid reference lengths are reported as unavailable, not silently omitted.
For the restricted ROS policy, success requires an explicit completion claim
verified using published target evidence and the benchmark distance criterion.

\paragraph{Required-interaction recall (ISR).}
Let $R_{ip}$ contain the required interaction IDs of valid plan $p$, and
$G_i$ contain the IDs with verified required effects. A failed macro can
contribute completed sub-effects; each ID contributes at most once per episode.
\begin{equation}
\mathrm{ISR}_i=\max_{p:\,R_{ip}\ne\emptyset}
 \frac{|G_i\cap R_{ip}|}{|R_{ip}|}.
\end{equation}
\begin{equation}
\mathrm{ISR}=\frac{1}{N_{\mathrm{req}}}
 \sum_{i:\,\mathrm{required}}\mathrm{ISR}_i.
\end{equation}
Episodes without required interactions are excluded from ISR. Opening effects
must satisfy the terminal opening threshold or a verified transient opening
condition. The complete-plan success indicator is reported separately.

\paragraph{Target-class interaction precision (IP).}
Let $A_i$ be the number of interaction attempts and $V_i$ the number producing
at least one verified physical effect in an annotated target class. Classes
combine channel/container domain and object category. An attempt contributes
at most once, including when a multi-joint macro is only partly successful.
An open effect crosses from below $0.8$ to at least $0.8$ during the attempt;
transient opening within a drawer scan also counts. Reopening after closure
can earn credit again; an effect-free repeated request cannot. Let $U_i=1$
only for episodes where interaction is unnecessary, and zero otherwise.
\begin{equation}
\mathrm{IP}_i=\begin{cases}V_i/A_i,&A_i>0,\\U_i,&A_i=0.\end{cases}
\qquad \mathrm{IP}=\frac{1}{N}\sum_i\mathrm{IP}_i.
\end{equation}
The object need not be the reference instance. IP jointly reflects target-class
selection and physical execution; successful exploration of other classes
receives no IP credit but is not automatically an execution error.

\paragraph{Normalized total cost.}
Let $E_i$ count attempts without any qualifying physical effect, including
failed and effect-free repeated requests. An attempt with a verified partial
effect is not an error, regardless of its category or macro acknowledgement.
Define the operation cost and normalized episode cost as
\begin{equation}
C_i^{\mathrm{op}}=L_i+\lambda A_i+\mu E_i,
\end{equation}
\begin{equation}
\mathrm{Cost}_i=\begin{cases}
\min(C_i^{\mathrm{op}}/B,1),&S_i=1,\\
1,&S_i=0.
\end{cases}
\end{equation}
\begin{equation}
\mathrm{TotalCost}=\frac{1}{N}\sum_i\mathrm{Cost}_i.
\end{equation}
The default protocol freezes $\lambda=0.3$, $\mu=1$, and $B=30$ before
evaluation. The budget is in weighted operation-cost units and is shared by
the compared methods. It is a scoring ceiling, not a rollout stopping rule;
SR and SPL retain their navigation-success definitions. Failures and successes
exceeding the budget have Cost one. We retain the unnormalized operation cost
and report SR alongside Cost. Metric versions and cost parameters are never
pooled across incompatible runs.
